\section{Results III: Training Convergence}
\label{sec:results-convergence}

Sections~\ref{sec:results-params} and~\ref{sec:results-timing} compare the models per evaluation; this section
checks that the hybrid models actually train, i.e.\ that they reach a loss and an accuracy comparable to the
classical PINN within a comparable training budget. It is a sanity check, not an accuracy benchmark: each
model is trained once (seed 0), on one PDE, with one training schedule.

\paragraph{Setup.}
All five models of Table~\ref{tab:configs} (the classical PINN and the four hybrid configurations of
Section~\ref{sec:configs}) are trained on all three boundary conditions and on both the forward and the inverse
problem (Section~\ref{sec:background}), six scenarios in total. Each scenario has a budget of $N_A$ Adam steps
and $N_L$ L-BFGS iterations: $N_A=5000$, $N_L=10000$ for Dirichlet and Neumann, and $N_A=N_L=20000$ with the
PDE residual weighted by 10 for Robin. \todo{reconcile with Section~\ref{sec:bg-loss}, which states that all
loss terms have unit weight.} The two model families spend this budget differently. The classical PINN runs
$N_A$ Adam steps at learning rate $10^{-3}$ and then L-BFGS (strong-Wolfe line search) for at most $N_L$
iterations. The hybrids skip L-BFGS, whose line search needs many extra circuit evaluations per iteration, and
instead run Adam for $N_A+N_L$ steps (learning rate $5\times10^{-3}$ for the circuit parameters and $10^{-3}$
for the rest) with cosine decay to $10^{-5}$. The optimiser therefore differs between the two families, and a
loss or accuracy gap between a hybrid and the classical PINN reflects both the model and its optimiser. To
separate the two, we also train the classical PINN with the hybrids' schedule: Adam for $N_A+N_L$ steps at
learning rate $10^{-3}$ with the same cosine decay, and no L-BFGS. We call this the \emph{classical (Adam)}
control; it shares the hybrids' optimiser and schedule, and differs from them in the network. Comparisons among
the four hybrids, which share the same optimiser, are free of the confound as well. For a given scenario every
model uses the same seed, so all of them see the same collocation points, initial data and
observations; only the network differs. Accuracy is measured against the reference solution as mean absolute
error (MAE).

\paragraph{Loss curves.}
Figure~\ref{fig:loss-curves} plots training loss (log scale) against loss evaluation for each scenario; for
the classical PINN's L-BFGS phase this counts line-search evaluations, not iterations. At the end of training
every hybrid configuration has a higher loss than the classical PINN in every scenario, by a factor of 1.6 to
57 for narrow, baseline, wide and deep alike, except narrow on the two Robin problems (270 and 390). Most of
this gap opens after the classical PINN switches to L-BFGS, where its loss drops by a further factor of 20 to
110 within 500 to 1700 evaluations. The classical (Adam) control, trained with the hybrids' optimiser,
shows no such drop. Against it, baseline and wide end at 0.2 times its loss on both Dirichlet problems and
within a factor of 0.9 to 2.2 on the Robin and inverse Neumann problems; only the forward Neumann problem
leaves them clearly above it (4.5 and 1.7 times). Narrow ends above the control in four of the six scenarios, by
up to a factor of 25 on Robin.

\paragraph{Accuracy.}
Table~\ref{tab:conv-mae} lists the test MAE. No hybrid configuration matches the classical PINN trained with
L-BFGS in any scenario: the best hybrid's MAE is 1.5 to 6.6 times the classical PINN's, and the worst's is
4.1 to 27 times. The optimiser accounts for most of this gap. Without L-BFGS, the classical PINN's own MAE
rises by a factor of 2.4 to 4.5, and against this classical (Adam) control the best hybrid is more accurate
in four of the six scenarios (0.63 to 0.95 times its MAE on the Dirichlet and Robin problems) and less
accurate only on the two Neumann problems (1.6 times). Baseline on its own beats the control in the same four
scenarios, with 59\% of its parameters (Table~\ref{tab:configs}). Narrow and deep do not beat it anywhere.
In this run, then, a hybrid of the right size is about as accurate as the larger classical network when both
use the same optimiser. The much larger gap to the Adam+L-BFGS column is what L-BFGS adds, and L-BFGS is the
optimiser the hybrids skip because of their per-evaluation cost (Section~\ref{sec:results-timing}).

Among the hybrids, no configuration is best everywhere. Baseline or wide is the most accurate in five of the
six scenarios, deep in the remaining one (Neumann, forward). Going from 2 to 4 qubits (narrow to baseline)
lowers the MAE in five of six scenarios, by up to a factor of 8 on the Robin problems; going from 4 to 6
qubits (baseline to wide) lowers it in three and raises it in three. Doubling the depth (baseline to deep)
raises the MAE in five of six scenarios. With a single seed these are observations about this run, not
trends; the per-seed spread of the hybrids' MAE is not known.

\begin{table}[t]
\centering
\caption{Test MAE after training (seed 0). Lower is better. \emph{Classical} uses Adam then L-BFGS; the
classical (Adam) control and all hybrids use Adam with cosine decay only. Bold marks the most accurate hybrid in
each row; an underline marks hybrids more accurate than the classical (Adam) control.}
\label{tab:conv-mae}
\small
\setlength{\tabcolsep}{4pt}
\begin{tabular}{@{}lrrrrrr@{}}
\toprule
 & \multicolumn{2}{c}{classical} & \multicolumn{4}{c}{hybrid} \\
\cmidrule(lr){2-3}\cmidrule(l){4-7}
Scenario & Adam+L-BFGS & Adam & narrow & baseline & wide & deep \\
\midrule
Dirichlet, fwd & $2.31\times10^{-4}$ & $5.43\times10^{-4}$ & $8.06\times10^{-4}$ & $\underline{\mathbf{3.40\times10^{-4}}}$ & $\underline{4.26\times10^{-4}}$ & $9.86\times10^{-4}$ \\
Neumann, fwd   & $1.66\times10^{-4}$ & $6.93\times10^{-4}$ & $1.85\times10^{-3}$ & $2.39\times10^{-3}$ & $1.31\times10^{-3}$ & $\mathbf{1.09\times10^{-3}}$ \\
Robin, fwd     & $8.96\times10^{-5}$ & $4.06\times10^{-4}$ & $2.03\times10^{-3}$ & $\underline{\mathbf{3.84\times10^{-4}}}$ & $4.88\times10^{-4}$ & $4.23\times10^{-4}$ \\
Dirichlet, inv & $2.87\times10^{-4}$ & $6.95\times10^{-4}$ & $7.13\times10^{-4}$ & $\underline{\mathbf{5.00\times10^{-4}}}$ & $\underline{6.76\times10^{-4}}$ & $1.16\times10^{-3}$ \\
Neumann, inv   & $2.49\times10^{-4}$ & $6.66\times10^{-4}$ & $2.80\times10^{-3}$ & $1.18\times10^{-3}$ & $\mathbf{1.05\times10^{-3}}$ & $1.70\times10^{-3}$ \\
Robin, inv     & $8.66\times10^{-5}$ & $3.42\times10^{-4}$ & $2.34\times10^{-3}$ & $\underline{2.90\times10^{-4}}$ & $\underline{\mathbf{2.72\times10^{-4}}}$ & $6.47\times10^{-4}$ \\
\bottomrule
\end{tabular}
\end{table}

\paragraph{Loss evaluations and wall-clock time.}
The classical PINN's L-BFGS phase meets its stopping tolerance well inside its budget in every scenario, after
500 to 1700 loss evaluations, so it evaluates the loss 5500 to 6200 times in total (Dirichlet and Neumann)
and about 21000 to 22000 times (Robin). Each hybrid evaluates it a fixed $N_A+N_L$ times: 15000 and 40000,
i.e.\ 1.8 to 2.7 times as often as the classical PINN.

Training a hybrid took from about 15 minutes (narrow, Dirichlet) to over 5 hours (wide, Robin) per scenario,
against at most 6 minutes for the classical PINN. We do not derive slowdown factors from these times. Each
comes from a single run on a workstation that was not reserved for the benchmark. Background load inflated
some of them: wide's inverse problems took 2.0 to 2.7 times as long as its forward problems with the same step
counts, against about 1.15 times for narrow. Section~\ref{sec:results-timing} gives the controlled
per-evaluation costs.

\begin{figure}[t]
\centering
\begin{subfigure}{0.32\linewidth}
  \includegraphics[width=\linewidth]{loss_dirichlet_fwd.png}
  \caption{Dirichlet, forward}
\end{subfigure}
\hfill
\begin{subfigure}{0.32\linewidth}
  \includegraphics[width=\linewidth]{loss_neumann_fwd.png}
  \caption{Neumann, forward}
\end{subfigure}
\hfill
\begin{subfigure}{0.32\linewidth}
  \includegraphics[width=\linewidth]{loss_robin_fwd.png}
  \caption{Robin, forward}
\end{subfigure}

\vspace{0.5em}
\begin{subfigure}{0.32\linewidth}
  \includegraphics[width=\linewidth]{loss_dirichlet_inv.png}
  \caption{Dirichlet, inverse}
\end{subfigure}
\hfill
\begin{subfigure}{0.32\linewidth}
  \includegraphics[width=\linewidth]{loss_neumann_inv.png}
  \caption{Neumann, inverse}
\end{subfigure}
\hfill
\begin{subfigure}{0.32\linewidth}
  \includegraphics[width=\linewidth]{loss_robin_inv.png}
  \caption{Robin, inverse}
\end{subfigure}
\caption{Training loss vs.\ loss evaluation (seed 0) for the classical PINN (Adam, then L-BFGS from the dotted
line), the classical (Adam) control and the four hybrid configurations (both Adam with cosine decay), for all
three boundary conditions and both
problem types. Each curve is a centred rolling median over 201 evaluations, which hides Adam's step-to-step
oscillation; the classical curve's Adam and L-BFGS phases are smoothed separately. The classical curve ends early
because L-BFGS meets its stopping tolerance.}
\label{fig:loss-curves}
\end{figure}
